\documentclass{article}
\usepackage[utf8]{inputenc}
\usepackage[T1]{fontenc}
\usepackage{lmodern}
\usepackage[polish]{babel}
\usepackage{amsmath}
\usepackage{tikz}
\usepackage{algorithm}
\usepackage{algpseudocode}
\usepackage{hyperref}
\usepackage{float}
\usepackage{graphicx}
\usepackage{mathtools}
\usepackage{amsmath}
\usepackage{amsfonts}
\usepackage{amsmath}
\usepackage{amsmath}
\usepackage{booktabs}
\usepackage[margin=1in]{geometry}
\usepackage{booktabs} % Dla ładniejszych linii w tabeli
\usepackage{array}




\title{Algorytmy metaheurystyczne - Lista 0}
\author{Wojciech Typer - 279730}
\date{}

\begin{document}
\maketitle

\subsection*{1000 losowych permutacji}
Na poniższych wykresach przedstawiono uzyskane rezultaty dla: Dżibuti, Kataru, Urugwaju, Sahary Zachodniej oraz zimbabwe. \\
Na osi X zaznaczono ilość grup, a na osi Y średnią z minimów uzyskanych dla konkretnych grup.

\vspace{2\baselineskip}

\begin{figure}[H]
    \centering
    
    \begin{minipage}{0.4\textwidth}
        \centering
        \includegraphics[width=\textwidth]{/home/wojteq18/Uni-6/Algorytmy_Metaheurystyczne/Lab/Lista_0/results/djibouti.png}
        \caption{Dżibuti}
        \label{fig:djibouti}
    \end{minipage}
    \hfill
    \begin{minipage}{0.4\textwidth}
        \centering
        \includegraphics[width=\textwidth]{/home/wojteq18/Uni-6/Algorytmy_Metaheurystyczne/Lab/Lista_0/results/qatar.png}
        \caption{Katar}
        \label{fig:qatar}
    \end{minipage}

    \vspace{0.5cm}

    \begin{minipage}{0.4\textwidth}
        \centering
        \includegraphics[width=\textwidth]{/home/wojteq18/Uni-6/Algorytmy_Metaheurystyczne/Lab/Lista_0/results/uruguay.png}
        \caption{Urugwaj}
        \label{fig:uruguay}
    \end{minipage}
    \hfill
    \begin{minipage}{0.4\textwidth}
        \centering
        \includegraphics[width=\textwidth]{/home/wojteq18/Uni-6/Algorytmy_Metaheurystyczne/Lab/Lista_0/results/western_sahara.png}
        \caption{Sahara Zachodnia}
        \label{fig:western_sahara}
    \end{minipage}

    \vspace{0.5cm}

    \begin{minipage}{0.4\textwidth}
        \centering
        \includegraphics[width=\textwidth]{/home/wojteq18/Uni-6/Algorytmy_Metaheurystyczne/Lab/Lista_0/results/zimbabwe.png}
        \caption{Zimbabwe}
        \label{fig:zimbabwe}
    \end{minipage}

    \caption{Zestawienie wyników dla wszystkich analizowanych państw}
    \label{fig:wszystkie_wyniki}
\end{figure}

\newpage
W uzyskanych wynikach czasem dochodzi do sytuacji w której drogi się krzyżują, co jest niepoządane. Taki problem można jednak rozwiązaćw następujący sposób: \\
\begin{figure}[H]
    \centering
    \begin{minipage}{0.48\textwidth}
        \centering
        \begin{tikzpicture}
            \draw[thick] (0,0) -- (4,4);
            \draw[thick] (0,4) -- (4,0);
            
            \filldraw[black] (0,0) circle (3pt) node[anchor=north east] {A};
            \filldraw[black] (4,4) circle (3pt) node[anchor=south west] {B};
            \filldraw[black] (0,4) circle (3pt) node[anchor=south east] {C};
            \filldraw[black] (4,0) circle (3pt) node[anchor=north west] {D};
            \filldraw[red] (2,2) circle (3pt) node[anchor=south] {X};
            
            \draw[->, thick] (0.5,0.5) -- (1.5,1.5);
            \draw[->, thick] (2.5,2.5) -- (3.5,3.5);
            \draw[->, thick] (0.5,3.5) -- (1.5,2.5);
            \draw[->, thick] (2.5,1.5) -- (3.5,0.5);
        \end{tikzpicture}
        \caption{Krzyżujące się krawędzie}
    \end{minipage}
    \hfill
    \begin{minipage}{0.48\textwidth}
        \centering
        \begin{tikzpicture}
            \draw[thick] (0,0) -- (0,4);
            \draw[thick] (4,0) -- (4,4);
            
            \draw[dashed, gray] (0,0) -- (4,4);
            \draw[dashed, gray] (0,4) -- (4,0);
            
            \filldraw[black] (0,0) circle (3pt) node[anchor=north east] {A};
            \filldraw[black] (4,4) circle (3pt) node[anchor=south west] {B};
            \filldraw[black] (0,4) circle (3pt) node[anchor=south east] {C};
            \filldraw[black] (4,0) circle (3pt) node[anchor=north west] {D};
            
            \draw[->, thick] (0,1.5) -- (0,2.5);
            \draw[->, thick] (4,1.5) -- (4,2.5);
        \end{tikzpicture}
        \caption{Połączenie optymalne}
    \end{minipage}
    \label{fig:2opt_improvement}
\end{figure}
Na rysunku 7 przedstawiono rozwiązanie nieoptymalne, w którym drogi A-B i C-D krzyżują się w punkcie X. Aby poprawić to rozwiązanie, można zastosować metodę 2-opt, która polega na zamianie końcówek dwóch krawędzi. W tym przypadku, zamiast łączyć A z B i C z D, można połączyć A z C oraz D z B. Dzięki temu drogi nie będą się krzyżować, co może prowadzić do lepszego rozwiązania problemu.

\vspace{2\baselineskip}

\subsection*{Cykl komiwojażera na podstawie MST}
Na poniższym wykresie przedstawiono wyniki znalezionego cyklu komiwojażera na podstawie MSTdla wszystkich analizowanych państw:
\begin{figure}[h!]
    \centering
    \includegraphics[width=0.7\textwidth]{/home/wojteq18/Uni-6/Algorytmy_Metaheurystyczne/Lab/Lista_0/results/ex4_results.png}
    \caption{Cykl komiwojażera na podstawie MST}
    \label{fig:obrazek1}
\end{figure}

\vspace{2\baselineskip}

Na podstawie uzyskanych wyników, można stwierdzić, że wynik uzykany na podstawie minimalnego drzewa rozpinającego jest znacznie lepszy od wyników uzyskanych dla losowych permutacji. 

\newpage
\subsection*{Wizualizacje}

\begin{figure}[H]
    \centering
    
    \begin{minipage}{0.4\textwidth}
        \centering
        \includegraphics[width=\textwidth]{/home/wojteq18/Uni-6/Algorytmy_Metaheurystyczne/Lab/Lista_0/results/viz/djiboutiv.png}
        \caption{Dżibuti}
        \label{fig:djibouti}
    \end{minipage}
    \hfill
    \begin{minipage}{0.4\textwidth}
        \centering
        \includegraphics[width=\textwidth]{/home/wojteq18/Uni-6/Algorytmy_Metaheurystyczne/Lab/Lista_0/results/viz/qatarv.png}
        \caption{Katar}
        \label{fig:qatar}
    \end{minipage}

    \vspace{0.5cm}

    \begin{minipage}{0.4\textwidth}
        \centering
        \includegraphics[width=\textwidth]{/home/wojteq18/Uni-6/Algorytmy_Metaheurystyczne/Lab/Lista_0/results/viz/uruguayv.png}
        \caption{Urugwaj}
        \label{fig:uruguay}
    \end{minipage}
    \hfill
    \begin{minipage}{0.4\textwidth}
        \centering
        \includegraphics[width=\textwidth]{/home/wojteq18/Uni-6/Algorytmy_Metaheurystyczne/Lab/Lista_0/results/viz/westernsaharav.png}
        \caption{Sahara Zachodnia}
        \label{fig:western_sahara}
    \end{minipage}

    \vspace{0.5cm}

    \begin{minipage}{0.4\textwidth}
        \centering
        \includegraphics[width=\textwidth]{/home/wojteq18/Uni-6/Algorytmy_Metaheurystyczne/Lab/Lista_0/results/viz/zimbabwev.png}
        \caption{Zimbabwe}
        \label{fig:zimbabwe}
    \end{minipage}

    \caption{Zestawienie wyników dla wszystkich analizowanych państw}
    \label{fig:wszystkie_wyniki}
\end{figure}


\end{document}